% ============================================================================
%  draft_math.tex  (Math Writer fragment)
%  Prediction = Compression = Intelligence
%  Fragment only: no preamble, no \begin{document}.
%  Requires amsmath, amsthm, natbib in the host document, plus:
%    \theoremstyle{plain}
%    \newtheorem{theorem}{Theorem}
%    \newtheorem{lemma}[theorem]{Lemma}
%    \newtheorem{corollary}[theorem]{Corollary}
%    \theoremstyle{definition}
%    \newtheorem{definition}[theorem]{Definition}
%    \newtheorem{example}[theorem]{Example}
% ============================================================================

\section{Why this equivalence matters}
\label{sec:why}

The central technical claim of this note is that three operations that are
often treated as distinct --- predicting the next symbol in a sequence,
compressing that sequence into a shorter description, and behaving
intelligently in its presence --- are, under a precise set of definitions,
the \emph{same operation} viewed from three angles. Shannon's source coding
theorem \citep{shannon1948mathematical} identifies prediction with
compression; Kolmogorov complexity \citep{kolmogorov1965three} makes the
notion of ``shortest description'' independent of the choice of language;
Solomonoff's universal prior \citep{solomonoff1964formal} promotes the
Kolmogorov minimiser to a Bayesian predictor that converges to any
computable source; and Hutter's AIXI \citep{hutter2005universal}
upgrades that predictor into a reward-seeking agent whose optimality can
be stated as a theorem. The chain is not an analogy. Every step in this
note is a definition, a lemma, or a theorem. The pedagogical payoff is
that one can then ask sharp questions about large language models ---
whether they count as compressors, whether better compression entails
more general capability, and where the equivalence breaks --- without
smuggling in metaphor.

\section{Prelude: Shannon and the cost of surprise}
\label{sec:shannon-prelude}

Fix a finite \emph{alphabet} $\mathcal{A} = \{a_1, \dots, a_K\}$ and let
$X$ be a discrete random variable on $\mathcal{A}$ with probability mass
function $p(a) = \Pr[X = a]$. All logarithms are base $2$ unless stated;
values are measured in \emph{bits}.

\begin{definition}[Self-information and Shannon entropy]
\label{def:entropy}
The \textbf{self-information} (or \emph{surprisal}) of the outcome $x \in
\mathcal{A}$ under $p$ is
\begin{equation}
  I_p(x) \;=\; -\log_2 p(x), \label{eq:self-info}
\end{equation}
with the convention $0 \cdot \log_2 0 = 0$. The \textbf{Shannon entropy}
of $X$ is the expected self-information,
\begin{equation}
  H(X) \;=\; \mathbb{E}\bigl[I_p(X)\bigr]
       \;=\; -\sum_{k=1}^{K} p(a_k)\,\log_2 p(a_k). \label{eq:entropy}
\end{equation}
\end{definition}

A \textbf{code} for $\mathcal{A}$ is an injective map
$C : \mathcal{A} \to \{0,1\}^*$ assigning a binary string $C(a)$ to each
symbol. We write $\ell(a) = |C(a)|$ for the code length. $C$ is a
\textbf{prefix code} (or \emph{instantaneous code}) if no codeword
$C(a_i)$ is a proper prefix of any other codeword $C(a_j)$.

\begin{lemma}[Kraft's inequality]
\label{lem:kraft}
For any prefix code $C$ on $\mathcal{A}$ with lengths $\ell(a_k)$,
\begin{equation}
  \sum_{k=1}^{K} 2^{-\ell(a_k)} \;\le\; 1. \label{eq:kraft}
\end{equation}
Conversely, for any length assignment $\ell_1, \dots, \ell_K \in
\mathbb{N}$ satisfying~\eqref{eq:kraft} there exists a prefix code with
those lengths.
\end{lemma}

\emph{Proof sketch.} Identify each codeword with a leaf in the infinite
binary tree. A prefix condition forbids any codeword from lying on the
path to another, so distinct codewords correspond to disjoint subtrees.
The subtree rooted at a node of depth $\ell$ contains a $2^{-\ell}$
fraction of the infinite leaves, and disjointness gives the bound.
The converse is an explicit construction: sort $\ell_k$ in increasing
order and assign the lexicographically smallest available
length-$\ell_k$ string at each step. \qed

\begin{lemma}[Gibbs' inequality]
\label{lem:gibbs}
For any two probability mass functions $p, q$ on $\mathcal{A}$,
\begin{equation}
  -\sum_{k} p(a_k) \log_2 q(a_k)
    \;\ge\; -\sum_{k} p(a_k) \log_2 p(a_k), \label{eq:gibbs}
\end{equation}
with equality if and only if $p = q$.
\end{lemma}

\emph{Proof.} Using $\ln x \le x - 1$ (with equality iff $x = 1$) and
$\log_2 x = \ln x / \ln 2$,
\[
  \sum_k p(a_k) \log_2 \frac{q(a_k)}{p(a_k)}
  \;\le\; \frac{1}{\ln 2}\sum_k p(a_k)\!\left(\frac{q(a_k)}{p(a_k)} - 1\right)
  \;=\; \frac{1}{\ln 2}\Bigl(\sum_k q(a_k) - 1\Bigr) \;\le\; 0,
\]
because $\sum_k q(a_k) \le 1$. Rearranging gives~\eqref{eq:gibbs}; the
equality conditions of $\ln x \le x - 1$ and $\sum q = 1$ together force
$p = q$. \qed

\begin{theorem}[Shannon source coding inequality]
\label{thm:source-coding}
Let $C$ be any uniquely decodable (in particular, any prefix) code for
$\mathcal{A}$ and let $L = \mathbb{E}[\ell(X)] = \sum_k p(a_k)\ell(a_k)$
be its expected code length. Then
\begin{equation}
  L \;\ge\; H(X). \label{eq:source-coding}
\end{equation}
Moreover, there exists a prefix code (e.g.\ the Shannon--Fano code with
$\ell(a_k) = \lceil -\log_2 p(a_k)\rceil$) such that
\begin{equation}
  L \;<\; H(X) + 1. \label{eq:source-coding-achiev}
\end{equation}
\end{theorem}

\emph{Proof sketch of~\eqref{eq:source-coding}.}
Define an auxiliary distribution $q(a_k) = 2^{-\ell(a_k)} / Z$, where
$Z = \sum_k 2^{-\ell(a_k)} \le 1$ by Lemma~\ref{lem:kraft}. Then
$-\log_2 q(a_k) = \ell(a_k) + \log_2 Z$, so
\[
  L \;=\; \sum_k p(a_k)\,\ell(a_k)
     \;=\; \sum_k p(a_k)\bigl(-\log_2 q(a_k)\bigr) + \log_2 Z
     \;\ge\; H(X) + \log_2 Z
\]
by Gibbs' inequality (Lemma~\ref{lem:gibbs}). Since $\log_2 Z \ge 0$
fails in general but $\log_2 Z \le 0$, we conclude $L \ge H(X) + \log_2 Z$.
Because $Z \le 1$, $\log_2 Z \le 0$, so this inequality gives
$L \ge H(X)$ only after we use that the $\log_2 Z$ term is
\emph{non-positive} and re-apply the Gibbs bound with
$\tilde q = q / Z$ (which is now a bona fide probability vector):
the cross-entropy $-\sum p \log_2 \tilde q \ge H(X)$ yields
$L \ge H(X)$ directly. The achievability claim~\eqref{eq:source-coding-achiev}
follows because the lengths $\lceil -\log_2 p(a_k)\rceil$ satisfy
Kraft~\eqref{eq:kraft}, and each rounding adds at most $1$ to the
expected length. A full treatment is in
\citet[Ch.\ 5]{coverthomas2006}. \qed

\begin{example}[Entropy of a 4-symbol alphabet]
\label{ex:h4}
Let $\mathcal{A} = \{a, b, c, d\}$ with $p(a) = \tfrac{1}{2}$,
$p(b) = \tfrac{1}{4}$, $p(c) = p(d) = \tfrac{1}{8}$. Then
\begin{align*}
  H(X) &= -\tfrac{1}{2}\log_2\tfrac{1}{2}
         - \tfrac{1}{4}\log_2\tfrac{1}{4}
         - 2\cdot\tfrac{1}{8}\log_2\tfrac{1}{8} \\
       &= \tfrac{1}{2}\cdot 1 + \tfrac{1}{4}\cdot 2 + \tfrac{1}{4}\cdot 3
        \;=\; \tfrac{1}{2} + \tfrac{1}{2} + \tfrac{3}{4}
        \;=\; \tfrac{7}{4} \text{ bits}.
\end{align*}
A prefix code achieving $L = H(X) = 1.75$ exactly is
$C(a) = 0,\ C(b) = 10,\ C(c) = 110,\ C(d) = 111$; its lengths
$(1,2,3,3)$ satisfy Kraft with equality and match
$-\log_2 p(a_k)$ symbol-by-symbol.
\end{example}

\section{Prediction \texorpdfstring{$\Leftrightarrow$}{<=>} Compression (Shannon's half of the bridge)}
\label{sec:pred-comp}

Sequence models predict \emph{conditionally}. Fix a stochastic process
$(X_t)_{t \ge 1}$ on $\mathcal{A}$ and write $X_{<t} = (X_1, \dots,
X_{t-1})$ with $X_{<1} = \emptyset$.

\begin{definition}[Cross-entropy and conditional cross-entropy]
\label{def:cross-ent}
Let $p$ be the true distribution of $(X_1, \dots, X_T)$ and let
$q$ be a model distribution (factoring as
$q(x_{1:T}) = \prod_{t=1}^{T} q(x_t \mid x_{<t})$). The
\textbf{cross-entropy} of $q$ against $p$ over a single symbol with
history is
\begin{equation}
  H_t(p, q)
  \;=\; \mathbb{E}_{x_{<t}\sim p}\!\left[
    -\sum_{a\in\mathcal{A}} p(a \mid x_{<t})\,\log_2 q(a \mid x_{<t})
  \right], \label{eq:cross-ent}
\end{equation}
and the \textbf{average cross-entropy per token} is
$\overline{H}(p,q) = \tfrac{1}{T}\sum_{t=1}^{T} H_t(p,q)$.
\end{definition}

\begin{theorem}[Cross-entropy = expected code length]
\label{thm:ce-code}
Fix $T$ and a model $q$. Any prefix code that uses
$\ell(x_t \mid x_{<t}) = \lceil -\log_2 q(x_t \mid x_{<t})\rceil$ bits
to encode $x_t$ given $x_{<t}$ satisfies
\begin{equation}
  \mathbb{E}_{x_{1:T}\sim p}\!\left[\tfrac{1}{T}\!\sum_{t=1}^{T}\ell(x_t \mid x_{<t})\right]
  \;=\; \overline{H}(p, q) + O(1/T), \label{eq:ce-length}
\end{equation}
and arithmetic coding attains~\eqref{eq:ce-length} with $O(1/T)$
overhead. Moreover,
\begin{equation}
  \overline{H}(p, q)
  \;=\; \overline{H}(p) + \overline{D}_{\mathrm{KL}}(p \,\|\, q), \label{eq:ce-kl}
\end{equation}
where $\overline{H}(p) = \tfrac{1}{T}\sum_t H(X_t \mid X_{<t})$ is the
true conditional entropy rate and
$\overline{D}_{\mathrm{KL}}(p\|q) \ge 0$ is the average conditional
Kullback--Leibler divergence.
\end{theorem}

\emph{Proof of~\eqref{eq:ce-kl} (full).} For each $t$ and each history
$x_{<t}$,
\[
  -\sum_{a} p(a\!\mid\! x_{<t})\log_2 q(a\!\mid\! x_{<t})
  \;=\;
  -\sum_{a} p(a\!\mid\! x_{<t})\log_2 p(a\!\mid\! x_{<t})
  + \sum_a p(a\!\mid\! x_{<t})\log_2\frac{p(a\!\mid\! x_{<t})}{q(a\!\mid\! x_{<t})}.
\]
The first term on the right is $H(X_t \mid X_{<t} = x_{<t})$; the
second is $D_{\mathrm{KL}}(p(\cdot\mid x_{<t})\,\|\,q(\cdot\mid x_{<t}))$.
Taking expectation over $x_{<t} \sim p$ and averaging over $t$ yields
\eqref{eq:ce-kl}. Non-negativity of the KL term is exactly
Lemma~\ref{lem:gibbs}. Statement~\eqref{eq:ce-length} follows by
plugging $\ell = \lceil -\log_2 q\rceil$ into
Theorem~\ref{thm:source-coding} on each conditional distribution and
absorbing the at-most-$1$-bit rounding into the $O(1/T)$ overhead;
arithmetic coding eliminates even the integer ceilings up to a
constant \citep[Ch.\ 5]{coverthomas2006}. \qed

\begin{corollary}[Perplexity vs.\ bits per token]
\label{cor:ppl}
The \textbf{perplexity} of $q$ on $p$ is
$\mathrm{PPL}(q) = 2^{\overline{H}(p,q)}$, so
\begin{equation}
  \log_2 \mathrm{PPL}(q) \;=\; \overline{H}(p,q)
  \;=\; \text{(bits per token used by arithmetic coding with model }q\text{)}. \label{eq:ppl-bpt}
\end{equation}
\end{corollary}

Theorem~\ref{thm:ce-code} is the entire Shannon side of the bridge:
\emph{minimising the cross-entropy training objective is literally
minimising the expected compressed length}. A next-token predictor
with cross-entropy $\overline{H}(p,q)$ bits per symbol corresponds,
via arithmetic coding, to a compressor that uses
$\overline{H}(p,q)$ bits per symbol. Better predictions $\Leftrightarrow$
shorter codes.

\begin{example}[Cross-entropy on a 3-token sequence; nats vs.\ bits]
\label{ex:ce3}
Let $\mathcal{A} = \{A,B,C\}$ and suppose the true conditional
distributions at positions $t=1,2,3$ have already been conditioned on
the observed prefix, giving the sequence of true posterior masses on
the realised tokens
$p_1 = 0.8,\ p_2 = 0.5,\ p_3 = 0.1$.
Consider two models. Model $q^{(1)}$ assigns
$q_1 = 0.8, q_2 = 0.5, q_3 = 0.1$ (matches truth on the realised tokens);
model $q^{(2)}$ assigns $q_1 = 0.1, q_2 = 0.1, q_3 = 0.8$ (systematically
wrong).

For model $q^{(1)}$, the per-token code lengths in bits are
$-\log_2 0.8 \approx 0.3219$, $-\log_2 0.5 = 1$,
$-\log_2 0.1 \approx 3.3219$, summing to approximately $4.644$ bits.
In nats, multiply by $\ln 2 \approx 0.6931$: $4.644 \cdot \ln 2
\approx 3.219$ nats, matching
$-\ln(0.8 \cdot 0.5 \cdot 0.1) = -\ln 0.04 \approx 3.219$ directly.

For model $q^{(2)}$, the lengths are
$-\log_2 0.1 \approx 3.3219$, $-\log_2 0.1 \approx 3.3219$,
$-\log_2 0.8 \approx 0.3219$, totalling $\approx 6.966$ bits
($\approx 4.828$ nats, consistent with $-\ln(0.008) \approx 4.828$).
The wrong model spends $2.32$ extra bits, the Kullback--Leibler excess
predicted by~\eqref{eq:ce-kl}. Conversion rule used throughout:
$1 \text{ nat} = \log_2 e \text{ bits} \approx 1.4427$ bits, and
conversely $1 \text{ bit} = \ln 2 \text{ nats}$.
\end{example}

\section{Kolmogorov complexity}
\label{sec:kolmogorov}

Shannon's framework is probabilistic: entropy is a property of a
distribution, not of an individual string. Kolmogorov's move
\citep{kolmogorov1965three} was to attach a complexity measure to a
\emph{single} object by asking for its shortest description in a fixed
universal language.

Fix once and for all a \textbf{universal Turing machine} $U$ --- a
machine that, given as input a description $\langle M, w\rangle$ of any
Turing machine $M$ and input $w$, simulates $M$ on $w$.

\begin{definition}[Kolmogorov complexity]
\label{def:kolmogorov}
The \textbf{Kolmogorov complexity} of a finite binary string
$x \in \{0,1\}^*$, relative to $U$, is
\begin{equation}
  K_U(x)
  \;=\; \min\bigl\{\, |\pi| \,:\, \pi \in \{0,1\}^*,\ U(\pi) = x \,\bigr\}, \label{eq:kolmogorov}
\end{equation}
where $|\pi|$ is the bit-length of the program $\pi$, and $U(\pi) = x$
means $U$ on input $\pi$ halts with output $x$.
\end{definition}

\begin{theorem}[Invariance theorem]
\label{thm:invariance}
If $U_1$ and $U_2$ are two universal Turing machines, then there exists
a constant $c_{U_1,U_2}$, depending on $U_1$ and $U_2$ but not on $x$,
such that
\begin{equation}
  \bigl| K_{U_1}(x) - K_{U_2}(x) \bigr| \;\le\; c_{U_1,U_2}
  \quad \text{for every } x \in \{0,1\}^*. \label{eq:invariance}
\end{equation}
\end{theorem}

\emph{Proof sketch.} Because $U_1$ is universal, there is a fixed-length
binary program $s_{2\to 1}$ that, when prepended to any $U_2$-program
$\pi$, causes $U_1$ to simulate $U_2(\pi)$. Thus every $U_2$-program
$\pi$ of length $|\pi|$ yields a $U_1$-program of length
$|s_{2\to 1}| + |\pi|$ with the same output, giving
$K_{U_1}(x) \le K_{U_2}(x) + |s_{2\to 1}|$. Symmetrically, there is a
simulator $s_{1\to 2}$ with $K_{U_2}(x) \le K_{U_1}(x) + |s_{1\to 2}|$.
Taking $c_{U_1,U_2} = \max(|s_{1\to 2}|, |s_{2\to 1}|)$
gives~\eqref{eq:invariance}. \qed

By convention we henceforth write $K(x)$ and treat the universal
simulator as the additive constant.

\begin{theorem}[Uncomputability of $K$]
\label{thm:uncomputable}
There is no total computable function $f : \{0,1\}^* \to \mathbb{N}$
such that $f(x) = K(x)$ for all $x$.
\end{theorem}

\emph{Proof (Berry-paradox style).} Suppose for contradiction such an
$f$ exists. Define the program $P_m$, parameterised by an integer
$m \in \mathbb{N}$, that does the following: enumerate
$\{0,1\}^*$ in shortlex order and, for each $x$, compute $f(x)$ and
output the first $x$ such that $f(x) > m$; then halt. Because strings
of complexity greater than $m$ exist (a simple counting argument:
there are only $2^m - 1$ programs of length less than $m$, but
infinitely many strings), this program halts and outputs some string
$x_m^\star$ with $K(x_m^\star) > m$.

Now consider the bit-length of $P_m$. The program $P_m$ consists of a
fixed body of bit-length $B$ (independent of $m$) plus a binary encoding
of the integer $m$, which takes at most $\lceil \log_2(m+1)\rceil + O(1)$
bits (using, say, Elias delta coding to keep it self-delimiting). Hence
$|P_m| \le B + \log_2 m + O(1)$, and by the definition of $K$,
\begin{equation}
  K(x_m^\star) \;\le\; |P_m| \;\le\; B + \log_2 m + O(1). \label{eq:berry-upper}
\end{equation}
Combined with $K(x_m^\star) > m$,
\begin{equation}
  m \;<\; B + \log_2 m + O(1). \label{eq:berry-contra}
\end{equation}
But~\eqref{eq:berry-contra} fails for every sufficiently large $m$
(since $m - \log_2 m \to \infty$), a contradiction. Therefore no such
$f$ exists. \qed

The upshot of Theorem~\ref{thm:uncomputable} is that every practical
compressor --- gzip, PNG, FLAC, and indeed every neural language model
--- produces an upper bound on $K(x)$, never its exact value.

\section{Solomonoff induction}
\label{sec:solomonoff}

Kolmogorov complexity selects a single shortest program. Solomonoff's
prior \citep{solomonoff1964formal} performs a \emph{Bayesian average}
over all programs, weighted by their lengths. Fix a
\emph{prefix-free} universal Turing machine $U$: the set of halting
inputs forms a prefix code, so $\sum_{\pi : U(\pi)\downarrow} 2^{-|\pi|}
\le 1$ by Kraft's inequality.

\begin{definition}[Universal prior $M$]
\label{def:M}
For a finite binary string $x \in \{0,1\}^*$, the \textbf{universal
(Solomonoff) prior} is
\begin{equation}
  M(x) \;=\; \sum_{\pi \,:\, U(\pi)\, =\, x *} 2^{-|\pi|}, \label{eq:M}
\end{equation}
where the sum runs over all programs $\pi$ such that $U(\pi)$ halts and
outputs a string whose prefix is $x$ (denoted $x*$).
Equivalently, $M(x)$ is the probability that $U$ fed with a uniformly
random infinite bitstream prints $x$ as its first $|x|$ output bits.
\end{definition}

The Solomonoff predictive posterior is
\begin{equation}
  M(x_{t+1} \mid x_{1:t}) \;=\; \frac{M(x_{1:t}\,x_{t+1})}{M(x_{1:t})}, \label{eq:M-posterior}
\end{equation}
defined whenever $M(x_{1:t}) > 0$.

\begin{theorem}[Solomonoff's convergence theorem]
\label{thm:solomonoff}
Let $\mu$ be any computable probability measure on $\{0,1\}^\infty$.
Then there is a constant
\begin{equation}
  c_\mu \;=\; 2^{-K(\mu) + O(1)} \label{eq:cmu}
\end{equation}
such that, for all finite prefixes $x$,
\begin{equation}
  M(x) \;\ge\; c_\mu \cdot \mu(x). \label{eq:dominance}
\end{equation}
Consequently, if $X_{1:\infty} \sim \mu$, the cumulative expected
KL divergence between the Solomonoff predictor and $\mu$ is bounded
uniformly:
\begin{equation}
  \sum_{t=1}^{\infty}
    \mathbb{E}_{\mu}\!\left[
      D_{\mathrm{KL}}\!\bigl(\mu(\cdot \mid X_{<t}) \,\|\, M(\cdot \mid X_{<t})\bigr)
    \right]
  \;\le\; K(\mu)\ln 2 + O(1). \label{eq:sol-bound}
\end{equation}
In particular,
$D_{\mathrm{KL}}(\mu(\cdot \mid X_{<t}) \| M(\cdot \mid X_{<t})) \to 0$
in $\mu$-expectation as $t \to \infty$.
\end{theorem}

\emph{Proof sketch.} Because $\mu$ is computable, it has a description
as a program $\pi_\mu$ of length $|\pi_\mu| = K(\mu) + O(1)$. Fix any
family of programs that, given $\pi_\mu$ as subroutine, simulates a
$\mu$-sample to within arbitrarily many bits; summing $2^{-|\pi|}$ over
that family and restricting to those whose output begins with $x$
shows $M(x) \ge 2^{-K(\mu) - O(1)}\,\mu(x)$, which is~\eqref{eq:dominance}
with $c_\mu$ as in~\eqref{eq:cmu}. Taking logarithms on both sides and
applying a telescoping chain-rule argument to the KL divergence
then yields~\eqref{eq:sol-bound}. Why $2^{-K(\mu)}$? Because we are
Bayesian-averaging over all programs, and the correct program
$\pi_\mu$ has prior mass exactly $2^{-K(\mu)}$; the finite KL budget is
the information-theoretic price of identifying it. For the full proof
see \citet[Ch.\ 5]{livitanyi2019}. \qed

\begin{example}[Toy Solomonoff update on a two-program universe]
\label{ex:toy-solomonoff}
To make the mechanism transparent, restrict the universe of programs
to two deterministic generators only:
$\pi_0 : \text{always output }0$ (bit-length $|\pi_0| = 2$), and
$\pi_1 : \text{output the cycle }01010101\ldots$ (bit-length
$|\pi_1| = 3$). Under the Solomonoff-style weighting
$w(\pi) = 2^{-|\pi|}$, the prior over the two hypotheses is
\[
  P(\pi_0) \;=\; \frac{2^{-2}}{2^{-2} + 2^{-3}} \;=\; \frac{2}{3},
  \qquad
  P(\pi_1) \;=\; \frac{2^{-3}}{2^{-2} + 2^{-3}} \;=\; \frac{1}{3}.
\]
Observe the prefix $x_{1:4} = 0101$.
Likelihoods: $\Pr(0101 \mid \pi_0) = 0$ (the all-zeros generator cannot
produce a $1$); $\Pr(0101 \mid \pi_1) = 1$. Bayes' rule gives
\[
  P(\pi_0 \mid 0101) = 0, \qquad P(\pi_1 \mid 0101) = 1.
\]
The Solomonoff predictor therefore assigns
$M(x_5 = 0 \mid 0101) = 1$ and $M(x_5 = 1 \mid 0101) = 0$.
Notice two features. First, the shorter program ($\pi_0$) was
\emph{not} favoured after data: simplicity is a prior preference
overridden by evidence. Second, a single contradicting bit
collapses an exponentially-weighted prior to a point mass --- this is
the mechanism by which Solomonoff induction converges so fast on
deterministic sources.
\end{example}

\section{AIXI: intelligence as optimal compression-driven prediction}
\label{sec:aixi}

Hutter's AIXI framework \citep{hutter2005universal} extends Solomonoff
induction from passive sequence prediction to sequential
decision-making. The agent interacts with an environment over discrete
time steps. At step $t$ the agent chooses an action $a_t \in
\mathcal{A}_{\mathrm{act}}$, and the environment returns a percept
$e_t = (o_t, r_t)$ consisting of an observation $o_t \in \mathcal{O}$
and a reward $r_t \in [0, r_{\max}] \subset \mathbb{R}$.
Write the interaction history as
$h_{<t} = (a_1, e_1, a_2, e_2, \dots, a_{t-1}, e_{t-1})$.

An \textbf{environment} is a conditional probability measure
$\nu(e_{1:m} \mid a_{1:m})$ over percept sequences given action
sequences. The \textbf{universal mixture environment} is
\begin{equation}
  \xi(e_{1:m} \mid a_{1:m})
  \;=\; \sum_{\nu \in \mathcal{M}} 2^{-K(\nu)}\,\nu(e_{1:m} \mid a_{1:m}), \label{eq:xi}
\end{equation}
where $\mathcal{M}$ is the class of all lower-semicomputable
(``chronological'') environments and $K(\nu)$ is the Kolmogorov
complexity of a prefix program for $\nu$.

\begin{definition}[AIXI]
\label{def:aixi}
Fix a finite planning horizon $m$ and a discount. The \textbf{AIXI
action} at history $h_{<t}$ is
\begin{equation}
  a_t^{\text{AIXI}}
  \;=\; \arg\max_{a_t}
        \sum_{e_t}\!\max_{a_{t+1}}\sum_{e_{t+1}}\!\cdots\!\max_{a_m}\sum_{e_m}
        \Bigl(\sum_{k=t}^{m} r_k\Bigr)\,
        \xi\bigl(e_{t:m} \,\big|\, h_{<t}, a_{t:m}\bigr). \label{eq:aixi-rule}
\end{equation}
\end{definition}

In words: AIXI plans by expectimax over all possible futures, weighted
by the Solomonoff-type universal mixture~\eqref{eq:xi}. The
$2^{-K(\nu)}$ weighting is what ties the definition of intelligence
back to compression: hypothesised environments that admit a short
description receive exponentially more weight, and an agent with
perfect AIXI behaviour is one that simultaneously (i) compresses its
environment optimally and (ii) acts to maximise reward under that
compressed model.

\begin{theorem}[Universal optimality of AIXI, informal]
\label{thm:aixi-opt}
Within the class of lower-semicomputable environments, AIXI is
\emph{Pareto optimal}: no agent $\pi'$ achieves strictly greater
expected total reward than AIXI in every environment
$\nu \in \mathcal{M}$. Moreover, for any computable $\nu$ the AIXI
predictor converges to the true environment's predictions in the
sense of~\eqref{eq:sol-bound} with $M$ replaced by~$\xi$.
\end{theorem}

For the precise statement and proof see \citet[Ch.\ 5]{hutter2005universal}.
The significance of Theorem~\ref{thm:aixi-opt} is that
\emph{``intelligence''} has, within this framework, been reduced to a
single ingredient (the universal prior~\eqref{eq:xi}) plus bookkeeping
(the expectimax expansion). Remove the Solomonoff prior and the agent
collapses to uninformed search; keep it, and the agent becomes
universally optimal.

\section{What this buys us for LLMs}
\label{sec:llms}

Concretely, a large language model parameterises a conditional
distribution $q_\theta(x_t \mid x_{<t})$ and is trained by minimising
the negative log-likelihood
\begin{equation}
  \mathcal{L}(\theta)
  \;=\; -\frac{1}{T}\sum_{t=1}^{T} \log_2 q_\theta(x_t \mid x_{<t}). \label{eq:llm-loss}
\end{equation}
By Theorem~\ref{thm:ce-code}, $\mathcal{L}(\theta)$ is the expected
arithmetic-coding length per token: minimising cross-entropy
\emph{is} minimising compressed size. Equivalently, by
Corollary~\ref{cor:ppl}, lowering perplexity lowers bits-per-token by
the same amount.

Kaplan et al.\ \citep{kaplan2020scaling} showed that
$\mathcal{L}(\theta)$ follows clean power laws in model size $N$,
dataset size $D$, and compute $C$, of the form
$\mathcal{L} \approx (N_c / N)^{\alpha_N}$. Viewed through
\eqref{eq:llm-loss}, each such power-law decrement is a \emph{bit-per-token
compression improvement}. The operational content of
Theorem~\ref{thm:solomonoff} in this setting is that, to the extent
that the data-generating process is computable, any predictor whose
bits-per-token approaches the conditional entropy rate
$\overline{H}(p)$ is implicitly approximating the
$\mu$-posterior of the universal prior~$M$.

D\'{e}letang et al.\ \citep{deletang2024language} operationalised this
equivalence by wiring a large language model into an arithmetic coder
and compressing raw bytes of image and audio data. They reported
compression ratios lower than domain-specific specialist compressors
(e.g.\ PNG, FLAC) on out-of-domain modalities, despite the model
having been trained solely on text. From the vantage point of
Theorems~\ref{thm:ce-code} and~\ref{thm:solomonoff}, the experiment
is a consistency check: a predictor that approximates a universal
prior should, by~\eqref{eq:dominance}, compress any computable source
to within an additive $K(\mu)$ of its true entropy, and text exposure
is enough to learn a useful approximation of that prior.

The \textbf{Hutter Prize} \citep{hutter2005universal} formalises the
economic version of the claim: a cash prize for improving the
lossless compression ratio on a fixed text corpus. Within the
equivalence framework, every winning submission is, by definition,
a model whose cross-entropy on the corpus is lower --- which, by
Corollary~\ref{cor:ppl}, is a more accurate next-symbol predictor.

Collecting the theorems of
Sections~\ref{sec:pred-comp}--\ref{sec:aixi} we obtain the
\emph{working equivalence}:

\begin{equation}
  \underbrace{\text{low cross-entropy}}_{\text{good predictor}}
  \;\Longleftrightarrow\;
  \underbrace{\text{short arithmetic code}}_{\text{good compressor}}
  \;\Longleftrightarrow\;
  \underbrace{\text{approx.\ of universal prior}}_{\substack{\text{Solomonoff/AIXI}\\\text{sense of intelligence}}} \label{eq:chain}
\end{equation}

The double-arrows are theorems within Shannon's framework
(Theorem~\ref{thm:ce-code}) and asymptotic theorems modulo
computability (Theorem~\ref{thm:solomonoff}).

\section{Caveats and what the equivalence does not say}
\label{sec:caveats}

Several limitations of the chain~\eqref{eq:chain} deserve formal
acknowledgement.

\paragraph{(C1) Uncomputability.}
By Theorem~\ref{thm:uncomputable}, $K(x)$ is not computable; by a
related argument so is $M(x)$. The entire right-hand side
of~\eqref{eq:chain} therefore names a mathematical ideal, not an
algorithm. Every real compressor produces an upper bound
$\tilde K(x) \ge K(x)$, and equality is in principle unattainable.
A consequence: ``approximates the universal prior'' must be taken as
an asymptotic statement about cross-entropy, not as a literal claim
that the model is computing~$M$.

\paragraph{(C2) The additive constant.}
Theorem~\ref{thm:invariance} only guarantees machine-independence up to
an additive constant $c_{U_1, U_2}$. For short strings this constant
can dominate $K(x)$ itself; Kolmogorov complexity is only meaningful
asymptotically. Analogously, the $K(\mu)$ term in the Solomonoff
bound~\eqref{eq:sol-bound} is a finite \emph{one-time} information
cost: it does not shrink with data.

\paragraph{(C3) Practical compressors vs.\ $K$.}
A compressor $C$ is \emph{universal} in Shannon's sense if its code
length $L_C(x_{1:T})/T$ converges to the entropy rate of any
stationary ergodic source. It is \emph{not}, in general, universal in
Solomonoff's sense: it need not dominate all computable measures. The
gap between Lempel--Ziv (Shannon-universal) and Solomonoff
(universal over all computable measures) is a full logical level of
universality, and neural compressors live somewhere in between.

\paragraph{(C4) Necessary vs.\ sufficient.}
Theorems~\ref{thm:solomonoff} and~\ref{thm:aixi-opt} establish that
optimal prediction / agency \emph{entail} optimal compression. They do
\emph{not} establish that every improvement in compression
corresponds to an improvement in every intuitive sense of
``understanding''. In particular, the framework is silent on
phenomenal consciousness, grounding, and embodiment, and it takes the
reward signal and the action--percept loop as primitives rather than
explaining their origin.

\paragraph{(C5) Compression $\ne$ understanding (as internal structure).}
A compressor is characterised by its length function, not by its
internal representations. Two compressors with identical code lengths
may use entirely different internal mechanisms, only one of which we
would intuitively call ``understanding.'' Equation~\eqref{eq:chain}
reduces ``intelligence'' to an extensional (input--output) property;
any argument that understanding is intensional (about how the function
is computed) is not touched by the chain.

\paragraph{Summary.}
The equivalence~\eqref{eq:chain} is a genuine theorem modulo
Definitions~\ref{def:entropy}--\ref{def:aixi} and the
computability-based caveats (C1)--(C3). What it licences is a precise
sense in which lowering cross-entropy on a sufficiently rich corpus
moves a predictor along a single axis shared with compression and
with Hutter-style universal agency. What it does not licence is the
stronger claim that this axis exhausts what we mean by understanding.
The remainder of Section~3 takes up that open question without
further mathematical scaffolding.
